% ============================================================================
% CHEATSHEET TEMPLATE
% ============================================================================
% Based on the ANA-G cheatsheet style
%
% Features:
% - 4-column landscape layout on A4
% - Colored boxes for main concepts and sub-concepts
% - Compact formatting optimized for information density
% - Mathematical typesetting support
%
% Usage:
% 1. Replace metadata (title, author, date)
% 2. Adjust number of columns if needed (see \Ncolumn)
% 3. Add your content using mainbox and subbox environments
% 4. Compile with pdflatex (run twice for proper layout)
% ============================================================================

\documentclass[a4paper, fontsize=7pt]{scrartcl}

% ============================================================================
% ENCODING & LANGUAGE
% ============================================================================
\usepackage[utf8]{inputenc}
% \usepackage[german]{babel}    % Uncomment for German hyphenation
% \usepackage[english]{babel}   % Uncomment for English hyphenation

% ============================================================================
% PAGE LAYOUT
% ============================================================================
% 4-column landscape layout with minimal margins
\usepackage[landscape, left=0.75cm, top=1cm, right=0.75cm, bottom=1.5cm, footskip=5pt]{geometry}
\usepackage{flowfram}
\ffvadjustfalse
\setlength{\columnsep}{1cm}

% Number of columns (adjust as needed: 3-4 recommended for A4)
\Ncolumn{4}

% Remove paragraph indentation
\setlength{\parindent}{0cm}

% ============================================================================
% COLORS & BOXES
% ============================================================================
\usepackage[most]{tcolorbox}

% Base style for all boxes
\tcbset{
  base/.style={
    boxrule=0.3mm,
    leftrule=1mm,
    left=1.75mm,
    arc=0mm,
    fonttitle=\bfseries,
    colbacktitle=black!10!white,
    coltitle=black,
    toptitle=0.75mm,
    bottomtitle=0.25mm,
    title={#1}
  }
}

% Main concept box (green accent)
% Use for: definitions, theorems, important formulas
\definecolor{maincolor}{rgb}{0.004, 0.5, 0.35}  % Nice green
\newtcolorbox{mainbox}[1]{
  colframe=maincolor,
  base={#1}
}

% Sub concept box (gray accent)
% Use for: examples, remarks, minor concepts
\newtcolorbox{subbox}[1]{
  colframe=black!20!white,
  base={#1}
}

% Alternative colors (uncomment to use)
% \definecolor{maincolor}{rgb}{0.2, 0.4, 0.8}   % Blue
% \definecolor{maincolor}{rgb}{0.8, 0.2, 0.2}   % Red
% \definecolor{maincolor}{rgb}{0.5, 0.2, 0.6}   % Purple

% ============================================================================
% MATHEMATICS
% ============================================================================
\usepackage{amsmath, amsthm, mathtools, amssymb}
\usepackage{amsfonts, amstext}
\allowdisplaybreaks  % Allow page breaks in equations

% Common math shortcuts (customize as needed)
\def\limn{\lim_{n\to \infty}}
\def\limxo{\lim_{x\to 0}}
\def\limxi{\lim_{x\to\infty}}
\def\limxn{\lim_{x\to-\infty}}
\def\sumk{\sum_{k=1}^\infty}
\def\sumn{\sum_{n=0}^\infty}

% Number sets
\def\N{\mathbb{N}}  % Natural numbers
\def\Z{\mathbb{Z}}  % Integers
\def\Q{\mathbb{Q}}  % Rationals
\def\R{\mathbb{R}}  % Reals
\def\C{\mathbb{C}}  % Complex

% Calculus
\def\dx{\,dx}
\def\dt{\,dt}
\def\dy{\,dy}
\newcommand{\ddx}[1]{\frac{d#1}{dx}}
\newcommand{\ddn}[2]{\frac{d#1}{d#2}}
\newcommand{\pdn}[2]{\frac{\partial #1}{\partial #2}}

% Big O notation
\def\oof{\mathcal{O}}

% ============================================================================
% TABLES
% ============================================================================
\usepackage{tabularx, multirow, booktabs}
\renewcommand*{\arraystretch}{1.5}  % More spacing in tables

% ============================================================================
% LISTS
% ============================================================================
\usepackage{enumitem}
\setitemize{itemsep=0.5pt, topsep=2pt}
\setenumerate{itemsep=0.75pt, topsep=2pt}

% ============================================================================
% GRAPHICS & TIKZ
% ============================================================================
\usepackage{graphicx}
\usepackage{tikz}
\usetikzlibrary{
  positioning,
  calc,
  fit,
  matrix,
  decorations,
  angles,
  quotes,
  babel,
  3d,
  arrows.meta,
  shadings,
  shapes.geometric
}

\tikzset{
  >=stealth,
  % Add custom TikZ styles here
}

% ============================================================================
% HYPERLINKS
% ============================================================================
\usepackage{hyperref}
\hypersetup{
  colorlinks=true,
  linkcolor=maincolor,
  urlcolor=maincolor
}

% ============================================================================
% ADDITIONAL PACKAGES
% ============================================================================
\usepackage{centernot}   % For \centernot command
\usepackage{marvosym}    % Additional symbols
\usepackage{wasysym}     % Additional symbols

% ============================================================================
% DOCUMENT METADATA
% ============================================================================
\title{Cheatsheet UGDM}
\author{Thomas Schnyder}
\date{September 2026}

% Table of contents depth
\setcounter{tocdepth}{4}
\setcounter{secnumdepth}{4}

% ============================================================================
% DOCUMENT START
% ============================================================================
\begin{document}

\maketitle

% ============================================================================
% SECTION 1: PERSÖNLICHKEIT UND FIRMENGRÜNDUNG
% ============================================================================
\section{Persönlichkeit und Firmengründung}

\begin{mainbox}{Unternehmerpersönlichkeit}
Typische Merkmale (Persönlichkeitstest "Bin ich ein Unternehmer-Typ?"):
\begin{itemize}
    \item Kontakteigenschaften, Begeisterungsfähigkeit
    \item Kontrolle (Impulse beherrschen), Risikobereitschaft
    \item Durchsetzungsvermögen, Selbstvertrauen/-sicherheit
    \item Selbstständigkeit, Schaffenskraft
    \item Führungseigenschaften, Neugierde
\end{itemize}
\textbf{Auswertung:} $\geq$10$\times$ Spalte 1/2 $\to$ ideale Unternehmerpersönlichkeit;
7--9$\times$ $\to$ Eigenschaften zu wenig ausgeprägt, intensive Arbeit nötig;
$<$6$\times$ $\to$ Eigenschaften fehlen, kaum entwickelbar.
\end{mainbox}

\begin{subbox}{Mariotti-Studie (ca. 400 Firmengründungen)}
Nützliche Charaktereigenschaften: Mut, Kreativität, Neugier, Entschlossenheit,
Disziplin, Flexibilität, Verantwortungsbewusstsein, Geduld, Enthusiasmus,
Einfühlungsvermögen.

\textbf{Am wichtigsten:} Fähigkeit, eine positive Grundeinstellung zu kultivieren.
\end{subbox}

\begin{subbox}{Goebel-Studie zu Jungunternehmern}
Erfolgreiche Jungunternehmer bringen eine Idee zum Reifen, indem sie beharrlich
und mit ungewöhnlicher Hartnäckigkeit immer wieder dasselbe Problem lösen wollen.
\end{subbox}

% ============================================================================
% SECTION 2: GRÜNDEN MIT KOMPONENTEN (I) – SMARTE GESCHÄFTSIDEEN
% ============================================================================
\section{Gründen mit Komponenten (I): Smarte Geschäftsideen}

\begin{mainbox}{Digitale Neukombination}
\textbf{Kernprinzip:} Digitalisierung ermöglicht, durch neuartige Kombination von
Bestehendem etwas Neues zu schaffen und am Markt zu verkaufen -- statt eine
komplett neue Erfindung zu machen.

\textbf{Typische Formen:}
\begin{itemize}
    \item Vorhandene Programmierleistung für neue, ursprünglich nicht vorgesehene
    Bereiche rekombinieren
    \item Physische Produkte mit digitalen Ressourcen koppeln
    \item Vorhandene Hardware + angepasste Software $\to$ höhere Leistungsfähigkeit
\end{itemize}
\end{mainbox}

\begin{mainbox}{Muster der Ideenentstehung (Studie Uni St. Gallen)}
Geschäftsideen unterscheiden sich oft nur gering von bekannten Geschäftsmodellen:
\begin{enumerate}
    \item \textbf{Übertragen:} alte Idee auf neuen Bereich (sehr häufig)
    \item \textbf{Kombinieren:} zwei alte Ideen neu kombiniert
    \item \textbf{Wiederholen:} erfolgreiches Vorgehen anderswo wiederholen
    \item Wirklich neue Ideen schaffen (sehr selten)
\end{enumerate}
\end{mainbox}

\begin{subbox}{Woraus ergibt sich die Neuartigkeit?}
Entweder aus dem Produkt/der Dienstleistung selbst, oder aus der \textbf{Art des
Verkaufs} (z.B. online statt offline). Das Zusammenwirken mehrerer Elemente
bietet das grösste Innovationspotenzial.
\end{subbox}

\begin{subbox}{Beispiele}
\textbf{Abo-Commerce (MyCouchbox):} Abo-Modell (Zeitschriftenwesen) + digitales
Kaufverhalten (E-Commerce) $\to$ monatliche Box mit wechselndem Inhalt.

\textbf{Teekampagne (Faltin):} statt vieler Sorten/Kleinpackungen/Zwischenhandel
$\to$ eine Sorte, Grosspackungen, Direktverkauf per Internet $\to$ tiefe Kosten,
hohe Margen.

\textbf{Hasenfarm:} übersättigter Markt (T-Shirts), eigene Marke, aber fremde
Technik (Webshop, Herstellung, Warenbewirtschaftung).
\end{subbox}

% ============================================================================
% SECTION 3: GRÜNDEN MIT KOMPONENTEN (II) – SMARTE LEISTUNGSERSTELLUNG
% ============================================================================
\section{Gründen mit Komponenten (II): Smarte Leistungserstellung}

\begin{mainbox}{Gründen in Komponenten}
\textbf{Definition:} Konzentration auf die eigentliche Geschäftsidee; für die
übrigen Arbeitsabläufe (Prozesse) werden Fähigkeiten/Techniken \textbf{anderer}
genutzt (Shop-Betreiber, Produktionsspezialist, Logistikanbieter,
Bürodienstleister, Einkauf über standardisierte Märkte/Börsen/Auftragsproduktion).

Traditionell bräuchte ein Gründer selbst Wissen in Rechnungswesen, Recht,
Personalführung, Marketing usw. -- smart gegründet wird dies fremdvergeben.
\end{mainbox}

\begin{mainbox}{Kostenvorteile}
\begin{itemize}
    \item Fast keine eigenen Investitionen nötig
    \item Fast keine Fixkosten
    \item Variable Kosten nur bei tatsächlichen eigenen Umsätzen
\end{itemize}
$\to$ hohe Kosteneinsparungen gegenüber etablierten Konkurrenten.
\end{mainbox}

\begin{mainbox}{Kostenbasis als Scheiterursache (ZHAW-Studie)}
\textbf{Wichtigster Grund} für das Scheitern von Start-ups: ungenügende Anpassung
der Kostenbasis (Fixkosten werden nicht schnell genug abgebaut).

\textbf{Kurzfristige Starrheit von Kosten:}
\begin{itemize}
    \item Rechtlich: z.B. Kündigungsfristen von Verträgen
    \item Betrieblich: Personalabbau stört eingespielte Teams, Demotivation,
    Reputationsverlust, Wiedereinstellungskosten
    \item Technisch: nötige Mindestkapazität
\end{itemize}

\textbf{Langfristige Starrheit -- spezifische Investitionen:} Investitionen in
Wissen, Informatik, Maschinen etc., die auf einen bestimmten Firmenzweck
ausgerichtet sind. Nicht firmenintern alternativ nutzbar, nicht firmenextern
verkaufbar $\to$ schränken Kostenanpassung stark ein.
\end{mainbox}

\begin{subbox}{Fixkosten vs. variable Kosten}
\renewcommand*{\arraystretch}{1.3}
\begin{center}
\begin{tabularx}{\linewidth}{lX}
\textbf{Fixkosten} & fallen unabhängig von Produktionsmenge an (Miete,
Kreditzinsen, Abschreibungen, Löhne) \\
\textbf{Variable Kosten} & steigen mit Produktionsmenge (Rohstoffverbrauch,
temporäre Arbeitskräfte) \\
\end{tabularx}
\end{center}
\textbf{Quasi-fixe Kosten} (z.B. Software-Personalkosten): formal variabel, aber
Fachpersonal lässt sich kaum kurzfristig ab-/aufbauen $\to$ verhalten sich wie
Fixkosten.
\end{subbox}

\begin{mainbox}{Ressourcenbegriff für Neukombination}
Sehr weit gefasst:
\begin{itemize}
    \item Elemente physischer Art (Rohstoffe, Materialien, Werkzeuge)
    \item Elemente der Biografie des Gründers (Erfahrungen, Qualifikationen,
    Fähigkeiten)
    \item Netzwerk des Entrepreneurs
    \item Ideen des Entrepreneurs
\end{itemize}
\end{mainbox}

\begin{subbox}{Beispiel: Yamo (natürliche Babynahrung)}
Auftragsproduktion (HPP-Anlage), Verpackungshersteller, Informatikdienstleister
(Webshop), Buchhaltungsdienstleister, Logistiker -- Gründer koordinieren und
kontrollieren nur.
\end{subbox}

% ============================================================================
% SECTION 4: SMARTE IDEENFINDUNG
% ============================================================================
\section{Smarte Ideenfindung}

\begin{mainbox}{Drei Ansatzpunkte der Ideenfindung}
\begin{enumerate}
    \item \textbf{Problemlösen:} Kundenperspektive einnehmen, Unzulänglichkeiten/
    ungedeckte Bedürfnisse erkennen
    \item \textbf{Trending:} Trends erkennen und nutzen (z.B. Ökologie,
    Individualität, Digitalisierung)
    \item \textbf{Neue Technologie/Erfindung kommerzialisieren} (oft zeit- und
    kostenintensiv)
\end{enumerate}
\end{mainbox}

\begin{subbox}{Ansatz Problemlösen -- 2 Schritte}
\begin{enumerate}
    \item Welche Probleme würde man gerne lösen?
    \item Welche Geschäftsideen würden diese Probleme beseitigen?
\end{enumerate}
\end{subbox}

\begin{subbox}{Merkregel}
Wirklich neue Ideen sind selten -- meist reicht es, eine bewährte Idee auf einen
neuen Bereich zu \textbf{übertragen} oder zwei bekannte Ideen zu
\textbf{kombinieren}. Der Erfolg entsteht oft durch \textbf{Neukombination
bestehender Komponenten}, nicht durch Neuerfindung.
\end{subbox}

% ============================================================================
% TIPS FOR AVOIDING OVERFULL BOXES
% ============================================================================
%
% 1. Keep formulas concise - split long equations over multiple lines
% 2. Use shorter variable names in inline math
% 3. Abbreviate table headers
% 4. Use align* environment for multi-line equations
% 5. Test compile frequently to catch overflow early
%
% Example of splitting a long formula:
% Instead of: $$a + b + c + d + e = f + g + h + i + j$$
% Use:
% \begin{align*}
% a + b + c + d + e &= f + g + h \\
%                   &\quad + i + j
% \end{align*}
%
% ============================================================================

\end{document}
